\documentclass[12pt]{article}
\usepackage{amsmath,amssymb,amsfonts}
\usepackage[margin=1in]{geometry}

\begin{document}

\begin{center}
{\Large\bfseries Chapter 10, Problem 12}
\end{center}

\noindent\textbf{Problem.} Show that the following three sets of functions, $\{\cos 3\theta\}$, $\{\sin 3\theta\}$ and $\{\cos 4\theta, \sin 4\theta\}$, transform among themselves under the elements of the group discussed in Problem 10.1. What are the representations and characters which are obtained from these three sets of functions? Are they irreducible representations? What can you say about the representation provided by $\{\cos 5\theta, \sin 5\theta\}$?

\bigskip
\noindent\textbf{Solution.}

The group from Problem 1 is the symmetry group of the equilateral triangle, $S_3 \cong D_3$, with 6 elements: $\{e, R, R^2, \sigma_a, \sigma_b, \sigma_c\}$, where $R$ is rotation by $120^\circ$ and $\sigma_a, \sigma_b, \sigma_c$ are reflections.

In terms of the angular coordinate $\theta$:
\begin{itemize}
\item $R$: $\theta \to \theta + 2\pi/3$
\item $R^2$: $\theta \to \theta + 4\pi/3$
\item $\sigma_a$: $\theta \to -\theta$ (reflection through line $a$ at $\theta = 0$)
\item $\sigma_b$: $\theta \to -\theta + 2\pi/3$
\item $\sigma_c$: $\theta \to -\theta + 4\pi/3$
\end{itemize}

\medskip
\noindent\textbf{Set 1: $\{\cos 3\theta\}$}

Under $R$: $\cos 3(\theta + 2\pi/3) = \cos(3\theta + 2\pi) = \cos 3\theta$.

Under $\sigma_a$: $\cos 3(-\theta) = \cos 3\theta$.

So $\cos 3\theta$ is invariant under all group operations. This gives a \textbf{one-dimensional representation} with all characters equal to 1:
\[
\chi(g) = 1 \quad \text{for all } g \in S_3.
\]
This is the trivial (identity) representation, $\Gamma_1$. It is \textbf{irreducible}.

\medskip
\noindent\textbf{Set 2: $\{\sin 3\theta\}$}

Under $R$: $\sin 3(\theta + 2\pi/3) = \sin(3\theta + 2\pi) = \sin 3\theta$.

Under $\sigma_a$: $\sin 3(-\theta) = -\sin 3\theta$.

Under $\sigma_b$: $\sin 3(-\theta + 2\pi/3) = \sin(-3\theta + 2\pi) = -\sin 3\theta$.

So rotations give $+1$ and reflections give $-1$. This is the \textbf{one-dimensional ``alternating'' representation} $\Gamma_2$:
\[
\chi(e) = 1, \quad \chi(R) = \chi(R^2) = 1, \quad \chi(\sigma_a) = \chi(\sigma_b) = \chi(\sigma_c) = -1.
\]
This is \textbf{irreducible}.

\medskip
\noindent\textbf{Set 3: $\{\cos 4\theta, \sin 4\theta\}$}

Under $R$ ($\theta \to \theta + 2\pi/3$):
\begin{align*}
\cos 4(\theta + 2\pi/3) &= \cos(4\theta + 8\pi/3) = \cos 4\theta \cos(8\pi/3) - \sin 4\theta \sin(8\pi/3) \\
&= \cos 4\theta \cdot(-\tfrac{1}{2}) - \sin 4\theta \cdot(-\tfrac{\sqrt{3}}{2}) \\
&= -\tfrac{1}{2}\cos 4\theta + \tfrac{\sqrt{3}}{2}\sin 4\theta.
\end{align*}
\begin{align*}
\sin 4(\theta + 2\pi/3) &= -\tfrac{\sqrt{3}}{2}\cos 4\theta - \tfrac{1}{2}\sin 4\theta.
\end{align*}

The representation matrix for $R$ is:
\[
D(R) = \begin{pmatrix} -1/2 & -\sqrt{3}/2 \\ \sqrt{3}/2 & -1/2 \end{pmatrix}, \quad \chi(R) = -1.
\]

Under $\sigma_a$ ($\theta \to -\theta$):
\[
\cos 4(-\theta) = \cos 4\theta, \quad \sin 4(-\theta) = -\sin 4\theta.
\]
\[
D(\sigma_a) = \begin{pmatrix} 1 & 0 \\ 0 & -1 \end{pmatrix}, \quad \chi(\sigma_a) = 0.
\]

The character table for this representation:
\[
\chi(e) = 2, \quad \chi(R) = \chi(R^2) = -1, \quad \chi(\sigma_a) = \chi(\sigma_b) = \chi(\sigma_c) = 0.
\]

This matches the two-dimensional irreducible representation of $S_3$. So this representation is \textbf{irreducible}.

\medskip
\noindent\textbf{The set $\{\cos 5\theta, \sin 5\theta\}$:}

Under $R$: $\cos 5(\theta + 2\pi/3) = \cos(5\theta + 10\pi/3)$. Since $10\pi/3 = 4\pi/3 \pmod{2\pi}$:
\[
D(R) = \begin{pmatrix} \cos(10\pi/3) & -\sin(10\pi/3) \\ \sin(10\pi/3) & \cos(10\pi/3) \end{pmatrix} = \begin{pmatrix} -1/2 & \sqrt{3}/2 \\ -\sqrt{3}/2 & -1/2 \end{pmatrix}.
\]
$\chi(R) = -1$.

Under $\sigma_a$: $\cos 5(-\theta) = \cos 5\theta$, $\sin 5(-\theta) = -\sin 5\theta$, so $\chi(\sigma_a) = 0$.

Characters: $\chi(e) = 2$, $\chi(R) = \chi(R^2) = -1$, $\chi(\sigma) = 0$.

This is the same set of characters as the two-dimensional irreducible representation. Therefore $\{\cos 5\theta, \sin 5\theta\}$ also provides the \textbf{irreducible two-dimensional representation} of $S_3$, equivalent to that provided by $\{\cos 4\theta, \sin 4\theta\}$.

\end{document}
